\subsection{An exceptional-prime finite pattern}
\begin{proposition}[Proved: modular affine solvability need not lift at one target]
\label{C4:exception}
There is a finite category with three objects, finite hom-sets and common
successors, and a triple of parallel arrows with an affine equalizer over
$\mathbb F_2$ but none over $\mathbb Q$. This assertion concerns the
specified family, not exactness of the category over $\mathbb F_2$.
\end{proposition}
\begin{proof}
Let the objects be $p,q,r$. In addition to identities, let the arrows
be $a_0,a_1,a_2:p\to q$, $b_1,\ldots,b_7:q\to r$, and
$c_0,\ldots,c_5:p\to r$. Define $b_j a_i=c_{e_{j,i}}$, where the
ordered triples $e_j$ are
\[
 012,\quad013,\quad015,\quad025,\quad124,\quad135,\quad234.
\]
These are all nonidentity compositions; identities provide the rest,
so associativity is immediate and is also checked exhaustively by the
program. Every object maps to $r$.

Over $\mathbb F_2$, take
$t=b_1+b_3+b_5+b_6+b_7$, of augmentation one. Its three marginals
coincide modulo two, as direct substitution shows.
For a rational certificate in the other direction, assign values on
$c_0,\ldots,c_5$ to two functionals:
\[
 \lambda_1=(-3/2,-1/2,-1/2,0,0,0),\qquad
 \lambda_2=(0,-1/2,0,0,1/2,0).
\]
Every listed triple $e_j=(x,y,z)$ satisfies
$\lambda_1(y)-\lambda_1(x)+\lambda_2(z)-\lambda_2(x)=1$.
Thus for $D(t)=(ta_1-ta_0,ta_2-ta_0)$,
$(\lambda_1,\lambda_2)D=\operatorname{aug}$ over $\mathbb Q$.
A rational kernel element has augmentation zero, so there is no
affine equalizer at $r$. At $q$ the only possible affine endomorphism
is the identity, which does not equalize the arrows; there are no
arrows $q\to p$. This proves the assertion at all targets.

The code \texttt{C4\_exception\_certificate.py} verifies the modular
vector, the rational functional, and every composable triple in this
19-arrow category. It also records why this is not a counterexample
to the central all-prime question: two distinct $b_j$ cannot be
equalized over any ring, since $r$ has only its identity endomorphism.
\end{proof}

\subsection{Contractible Kan simplex categories are Ramsey}
\begin{theorem}[Proved: the contractible Kan candidate is already exact]
\label{C4:contractible-kan-ramsey}
Let $K$ be a nonempty contractible Kan complex. Then the category
$\Delta_{\mathrm{inj}}/K$ has the morphism Ramsey property. Consequently
its colimit functor is exact over every unital ring. In particular,
choosing a contractible Dold--Kan complex from an acyclic but nonsplit
chain complex does not give a total nonexact example in this family.
\end{theorem}
\begin{proof}
Fix objects $A,B$ and a finite number $c$ of colours. If there is no
arrow $A\to B$, the conclusion is immediate with target $B$. Otherwise
let $L$ be the finite relational language with one relation $R_\sigma$
for each simplex label occurring on a nonempty face of $A$ or $B$;
its arity is the number of vertices of that face. Equal labels give
one symbol. This set of symbols is closed under all face maps.

Consider finite $L$-structures with the following rules. Every vertex
has exactly one vertex label. Every relation tuple has distinct
entries; it has at most one label of its arity. If an ordered tuple
has label $\sigma$, every ordered subtuple has the corresponding
face label of $\sigma$. No global order on these tuples is required
at this stage, and different orderings of the same vertex set may
be labelled independently, subject only to their common ordered
subtuples. This is a hereditary free amalgamation class: in an
amalgam all relation tuples come from one factor, and their subtuples
come from that factor as well; labels on a tuple in both factors
already agree on the common induced substructure. There are no new
cross-factor relation tuples.

Add an arbitrary linear order. Theorem~1.3 of
Evans--Hubi\v cka--Ne\v set\v ril \cite{D-EHN2021} gives a finite
ordered member $C$ with $C\longrightarrow(B)^A_c$, where $A$ and $B$
are encoded by labelling precisely their increasing tuples. Remove
from $C$ every relation tuple that is not increasing in its added
linear order. The remaining labels give a downward closed simplicial
subcomplex $L_C$ of the full simplex on $V(C)$ and a coherent
simplicial map
\[
 L_C\longrightarrow K.
\]
Deleting the nonincreasing tuples does not change any existing
induced copy of $A$ or $B$, since these structures have no such
tuples. In particular the specific copies used in the Ramsey
conclusion remain available.

The map $K\to *$ is a trivial Kan fibration: it is a Kan fibration
and a weak equivalence. Thus it has the right lifting property for
all simplicial monomorphisms, by
\citet[Theorem I.7.10, pp.~32--35]{ABGoerssJardine} and the
boundary-by-boundary extension described in its proof. Apply this to
$L_C\hookrightarrow\Delta^{V(C)}$ to extend the displayed map to the
full simplex. The extension defines an object $D$ of
$\Delta_{\mathrm{inj}}/K$ on the same ordered vertex set. Every
previous copy of $A$ or $B$ remains an induced copy, because all of
its faces were already labelled and the extension preserves those
labels.

Now colour the arrows $A\to D$ with $c$ colours. Restrict the
colouring to the old copies $A\to C$. The Ramsey property of $C$
gives an old copy $B\to C$ whose $A$-subcopies are monochromatic.
This is a copy $B\to D$ with the same property. Hence $D$ is the
required Ramsey target. The category has finite hom-sets and joint
embeddings: the same extension argument applies to a disjoint union
of two labelled simplices. The integral affine-equalizer proof of
Ramsey exactness applies, and then works over every unital ring.
\end{proof}

\subsection{A counterexample to the all-prime implication}

\begin{lemma}[Proved: finite marginal lifting in a perfect group algebra]
\label{C4:perfect-marginals}
Let $p$ be a prime, let $V$ be an abelian group on which multiplication
by $p$ is bijective, and put $k=\mathbb F_p$. Suppose
$\rho_i:V\twoheadrightarrow V_i$ is a finite family of surjective
homomorphisms, with each $V_i$ uniquely $p$-divisible. Let
$J_i=\ker(k[V]\to k[V_i])$. Then the finite lattice of ideals generated
from the $J_i$ by sums and intersections is distributive. Consequently,
a family $z_i\in k[V_i]$ lifts to one $z\in k[V]$ if its images agree
in $k[V]/(J_i+J_j)$ for every pair $i,j$. The lift is a finite linear
combination of elements of $V$.

The same statement holds for affine surjections of torsors under these
groups: choose any origin in the source torsor, and the induced origins
in its quotients.
\end{lemma}
\begin{proof}
The Frobenius map on $k[V]$ sends $\sum a_v[v]$ to
$\sum a_v[pv]$. It is bijective, with inverse obtained by replacing
$v$ by its unique $p$th root. The same is true in every $k[V_i]$.
Thus each kernel $J_i$ is stable under Frobenius and its inverse.
Sums and intersections of ideals with this property retain it:
for sums use $(x+y)^{1/p}=x^{1/p}+y^{1/p}$, and for intersections
apply the inverse in each ideal. Every such ideal is radical. Indeed,
if $x^n$ belongs to it, then so does $x^{p^m}$ for $p^m\geq n$;
applying the inverse Frobenius $m$ times puts $x$ in the ideal.

The radical ideals of a commutative ring form a distributive lattice
with intersection as meet and radical of the sum as join. Here all
sums under consideration are already radical, so the operations are
ordinary sum and intersection. Equivalently, one can verify
\[
 (I+J)\cap(I+J')=I+(J\cap J')
\]
by taking radicals and using the corresponding identity of closed
sets in the prime spectrum. This proves the claimed distributivity.

For the lifting assertion induct on the number of ideals. Suppose
$x$ lifts the first $m-1$ residues and $y$ lifts the last. Pairwise
compatibility gives $x-y\in J_i+J_m$ for every $i<m$, whence
\[
 x-y\in\bigcap_{i<m}(J_i+J_m)
      =\left(\bigcap_{i<m}J_i\right)+J_m.
\]
Subtracting the first summand from $x$ makes it agree with $y$ modulo
$J_m$ without changing any earlier residue. All elements of a group
algebra have finite support. Changing an affine origin only translates
a basis, so the torsor formulation follows.
\end{proof}

\begin{lemma}[Proved: coherent modular weights on rational primitives]
\label{C4:primitive-weights}
Fix $r\geq2$ and a prime $p$. For every finite ordered rational
$r$-cocycle object $S$ there is a finitely supported element
\[
 \mu_S\in\mathbb F_p[\operatorname{Prim}(S)],\qquad
 \operatorname{Prim}(S)=
 \{b\in C^{r-1}(\Delta^S;\mathbb Q):db=c_S\},
\]
of augmentation one. These elements may be chosen simultaneously
so that every induced embedding $T\to S$ pushes $\mu_S$ to $\mu_T$
by restriction of primitives.
\end{lemma}
\begin{proof}
The primitive set is a nonempty affine torsor under the rational
vector space $V_S=Z^{r-1}(\Delta^S;\mathbb Q)$. Include the empty
structure for this construction only, with its unique primitive.
Restriction to a subset is surjective: a difference of two primitives
on that subset is a positive-degree cocycle, hence a coboundary;
extend its primitive cochain to $S$ and correct a chosen primitive
on $S$.

We need a precise two-subset assertion. If $T,U\subseteq S$, any
primitives on $T$ and $U$ agreeing on $T\cap U$ extend to a primitive
on $S$. Subtract the restriction of a fixed primitive on $S$. The
result is an $(r-1)$-cocycle on
$\Delta^T\cup\Delta^U$. That complex is contractible when its
intersection is nonempty and is a disjoint union of at most two
simplices otherwise. Its cohomology in degree $r-1>0$ vanishes.
Write the difference as a coboundary, extend the preceding cochain
to the full simplex, and add its coboundary to the fixed primitive.
This also proves, on the translation spaces,
\[
 \ker(V_S\to V_T)+\ker(V_S\to V_U)
   =\ker(V_S\to V_{T\cap U}).                       \tag{C4.1}
\]
For example, to prove the nontrivial inclusion, take $v$ vanishing
on the intersection and extend the compatible pair $(v|_T,0|_U)$
to $w\in V_S$; then $v=(v-w)+w$ is the required decomposition.
The kernels in the group algebra therefore satisfy
$J_T+J_U=J_{T\cap U}$: the kernel of a group-algebra quotient is
generated by $[h]-1$ for elements $h$ of its group kernel.

Choose the weights by induction on the size of $S$, one choice for
each ordered isomorphism type. At size zero choose the unique basis
element. Suppose the choices have been made naturally on all smaller
structures. For a structure $S$, prescribe the previously chosen
weight on each proper subset. These marginals agree after restriction
to every intersection, by the inductive hypothesis. Equation (C4.1)
and Lemma~\ref{C4:perfect-marginals} give an element of
$\mathbb F_p[\operatorname{Prim}(S)]$ with all the prescribed
marginals. Include the empty marginal, so its augmentation is one.
A finite linearly ordered structure has no nonidentity increasing
self-isomorphism. Thus transporting this choice from the representative
to every ordered copy is unambiguous, and supplies all embedding
compatibilities. This completes the induction. The argument gives
finite support at every individual object, although no uniform bound
on support sizes is asserted.
\end{proof}

\begin{theorem}[Refuted: all-prime exactness forces rational exactness]
\label{C4:counterexample}
For each $r\geq2$, the category $A=\mathcal A_r(\mathbb Q)$ of finite
nonempty ordered rational $r$-cocycles is exact over $\mathbb F_p$
for every prime $p$, but is not exact over $\mathbb Q$ or $\mathbb Z$.
It has finite hom-sets and common successors. It lies in setting (S)
and satisfies (Rig), (SA), (FA), (St), (Op), (WEI), and (Fin).
It is not total, since $BA\simeq K(\mathbb Q,r)$.

Consequently the main all-prime implication in C.4 is false, both
with and without the finite-hom hypothesis, and remains false in
setting (S) with (Rig). Property (P) does not imply totality, even
for these nonempty connected finite-hom rigid strong Fra\"iss\'e
categories. The version whose index category is assumed total is
not decided by this example.
\end{theorem}
\begin{proof}
Fix $p$. The expansion category $B$ consists of finite ordered
$(r-1)$-cochains with rational values. Its reduct sends $b$ to $db$.
For each finite collection of objects and arrows, only finitely many
cochain values occur. Encode these values by finitely many relation
symbols, with value zero represented by absence of a relation.
The corresponding class of coloured $r$-uniform hypergraphs has free
amalgamation, so its arbitrary linear-order expansion is Ramsey by
\cite[Theorem~1.3]{D-EHN2021}. Hence every finite family of parallel
arrows in $B$ has an integral affine equalizer, by the rechecked
Ramsey-to-affine argument. Enlarging the finite set of colours to
include zero preserves all required objects and embeddings.

We give the finite-family passage explicitly, so no inference from
constant-coefficient or action-component homology is needed. Let
$a_0,\ldots,a_m:P\to Q$ be any finite parallel family in $A$.
Write the finite support of Lemma~\ref{C4:primitive-weights} as
\[
 \mu_Q=\sum_{b\in\mathcal S}w_b[b],\qquad
 \sum_b w_b=1\quad\text{in }\mathbb F_p.
\]
Jointly embed the finitely many primitive expansions $Q_b=(Q,b)$
into one cochain object $X^*$, by maps $j_b:Q_b\to X^*$.
For every primitive expansion $P_d$ occurring as some
$a_i^*b$, consider all the maps
\[
 j_b a_i:P_d\longrightarrow X^*
       \quad\text{with }a_i^*b=d.
\]
There are finitely many such types $d$ and finitely many such maps.
Use integral affine equalizers in $B$ successively for these families,
postcomposing later families with the previously obtained affine
arrows. Equivalently, at a later stage equalize the finitely many
ordinary arrows in the supports of the corresponding affine arrows.
This yields one affine arrow $s:X^*\to Y^*$ over $\mathbb F_p$
for which all the displayed maps of each fixed type $d$ have the
same image $v_d$. Postcomposition preserves every equality already
achieved.

Let $u:B\to A$ forget the primitive and set
\[
 t=\sum_{b\in\mathcal S}w_b\,u(s)u(j_b)
      \in\mathbb F_p[A(Q,uY^*)].
\]
Its augmentation is one. For every $i$ its composite with $a_i$ is
\[
 ta_i=\sum_d\left(\sum_{b:a_i^*b=d}w_b\right)u(v_d)
      =\sum_d\mu_P(d)\,u(v_d),
\]
independent of $i$, by coherence of the weights. Thus every finite
parallel family has an affine equalizer. The affine criterion proves
exactness over $\mathbb F_p$ for arbitrary module diagrams.

For rational nonexactness one can use a finite named family, not just
the nonzero homology of the base. On $Q=\{0,\ldots,2r+1\}$ choose an
$(r-1)$-cochain taking value $1$ on $(1,\ldots,r)$ and zero on every
other increasing $r$-tuple, and give $Q$ its coboundary cocycle.
The first $(r+1)$-point block has cocycle value one; the disjoint
second block has value zero. Let $P$ have $r$ points and zero
$r$-cocycle, and take the $2(r+1)$ embeddings $P\to Q$ onto the
codimension-one faces of these two blocks. At any target $D$, choose
a rational primitive $b_D$. If an affine combination $t$ equalized
all these face embeddings, evaluating $b_D$ and alternating over
the faces of the two blocks would give equal answers. The coboundary
identity instead gives $\operatorname{aug}(t)=1$ on the first block
and zero on the second. This contradiction excludes a rational
affine equalizer at every target. Clearing denominators also excludes
integral exactness.

The strong-class properties and the identification
$BA\simeq K(\mathbb Q,r)$ are the audited constructions of R13:
all subsets are induced strong substructures, cochains extend across
two simplices, nonempty-base amalgamation is functorial, and the
three-piece extension uses $H^r(S^1;\mathbb Q)=0$ for $r\geq2$.
The minimal finite-support argument gives (Op), and rigidity comes
from the order. For (Fin), if $G_C\leq G_B$ has finite index, every
point of $C$ has finite $G_B$-orbit. Free duplication over $B$ forces
$C\subseteq B$, hence $G_B\leq G_C$ and equality. (WEI) follows
from (Op). Finite hom-sets are subsets of sets of increasing
injections between finite ordered sets. Joint embeddings give
common successors. Finally $H_r(BA;\mathbb Z)\cong\mathbb Q\ne0$,
so $BA$ is not contractible and $A$ is not total. The established
profinite/all-prime equivalence gives (P).
\end{proof}

\begin{corollary}[Proved: the complete field profile of cocycle classes]
\label{C4:field-profile}
Let $\Lambda$ be a countable abelian group and let $r\geq2$. Then
\[
 \mathcal A_r(\Lambda)\text{ is exact over }\mathbb F_p
 \quad\Longleftrightarrow\quad
 p:\Lambda\longrightarrow\Lambda\text{ is bijective},
\]
and
\[
 \mathcal A_r(\Lambda)\text{ is exact over }\mathbb Q
 \quad\Longleftrightarrow\quad \Lambda\text{ is torsion}.
\]
Thus it is exact over every prime field $\mathbb F_p$ precisely when
$\Lambda$ is a rational vector space, and it is integrally exact
precisely when $\Lambda=0$. For a nonzero rational vector space these
categories give the counterexamples in Theorem~\ref{C4:counterexample}.
\end{corollary}
\begin{proof}
If multiplication by $p$ is bijective on $\Lambda$, the primitive
translation group
$Z^{r-1}(\Delta^S;\Lambda)$ is uniquely $p$-divisible. This follows
either from the cone contraction, which identifies it with a finite
power of $\Lambda$, or by applying the inverse of multiplication by
$p$ to cochains coordinatewise. The two-subset extension and the
perfect group-algebra argument of
Lemmas~\ref{C4:perfect-marginals}--\ref{C4:primitive-weights} therefore
apply verbatim. The finite-family proof then gives exactness over
$\mathbb F_p$ for every module diagram.

For necessity, exactness annihilates all positive derived colimits,
in particular homology with the constant $\mathbb F_p$ diagram.
The base has $H_r(BA;\mathbb Z)=\Lambda$, and no reduced integral
homology below degree $r$. The universal coefficient theorem gives
\[
 H_r(BA;\mathbb F_p)=\Lambda/p\Lambda,
 \qquad
 H_{r+1}(BA;\mathbb F_p)\twoheadrightarrow\Lambda[p].
\]
Both groups on the right must vanish. Thus multiplication by $p$ is
surjective and injective. This use of constant coefficients proves
necessity only; sufficiency was established above for arbitrary
modules.

If $\Lambda$ is torsion, the finitely many values occurring in a
fixed finite parallel family generate a finite subgroup
$\Lambda_0$. The finite-group primitive-averaging proof in
Theorem~\ref{A:finite-cocycles} supplies a rational affine equalizer
inside $\mathcal A_r(\Lambda_0)$, hence in $\mathcal A_r(\Lambda)$.
The affine criterion proves rational exactness. Conversely, if
$\Lambda$ is not torsion, choose a nonzero linear functional on
$\mathbb Q\otimes\Lambda$ and normalize it at a nonzero image of
some $\lambda\in\Lambda$. It induces a homomorphism
$\chi:\Lambda\to\mathbb Q$ with $\chi(\lambda)=1$.
Apply $\chi$ to a target primitive in the finite two-block face
obstruction of Theorem~\ref{C4:counterexample}. The same alternating
sum yields $1=0$, so rational exactness fails.

Multiplication by every prime is bijective exactly for rational
vector spaces. Integral exactness would force the constant homology
$H_r(BA;\mathbb Z)=\Lambda$ to vanish; for $\Lambda=0$ the category
is $\Delta_{\mathrm{inj}}$, which is integrally exact. This proves
the remaining assertions.
\end{proof}

\begin{example}[A finite modular primitive distribution]
\label{C4:triangle-distribution}
For a triangle of cocycle value one, write a primitive as
$(x,y,z)=(b_{01},b_{12},b_{02})$, so $x+y-z=1$.
For any prime $p$, put weight $+1$ on
\[
 (1/p,j/p,(j+1-p)/p),\qquad 0\leq j<p,
\]
and weight $-1$ on
\[
 (0,j/p,(j-p)/p),\qquad 1\leq j<p.
\]
This has $2p-1$ distinct atoms and integral augmentation one. Its
three marginals are respectively
$p[1/p]-(p-1)[0]$, $[0]$, and $[0]$.
Modulo $p$ all are $[0]$. This exhibits explicitly how rational
denominators create the modular marginal cancellation, while the
rational cocycle still has value one. The exact-arithmetic program
\texttt{C4\_perfect\_marginal\_certificate.py} checks the formulas
for eight primes; the displayed calculation proves them uniformly.
\end{example}

\begin{proposition}[Proved: a bounded-denominator criterion]
\label{C4:noetherian}
Let $A$ be nonempty with common successors, and fix a finite parallel
family $a_0,\ldots,a_k:p\to q$. Write
\[
 D:\mathbb Z[A(q,-)]\longrightarrow
       \mathbb Z[A(p,-)]^k,
 \qquad D(t)=(ta_i-ta_0)_{1\leq i\leq k},
\]
and let $S$ be the pointwise saturation of $\operatorname{im}D$ in
the target. If $S/\operatorname{im}D$ is annihilated by a positive
integer $N$, an affine equalizer over $\mathbb F_p$ for any prime
$p\nmid N$ implies a rational affine equalizer for this family.
Consequently all-prime exactness implies rational exactness whenever
these torsion quotients have bounded exponent for every finite family.

The bounded-exponent condition holds if each such saturation functor
is generated by finitely many elements at finitely many objects;
in particular it holds if finite sums of integral representable
functors are Noetherian. In setting (S), an alternative sufficient
condition is that every finite sum of the smooth permutation modules
$\mathbb Z[G/G_P]$ is Noetherian as a $\mathbb Z[G]$-module.
No Noetherian hypothesis is asserted for arbitrary oligomorphic
groups or for the independently ordered vector-space example.
\end{proposition}
\begin{proof}
Choose a modular affine equalizer at a target $r$ and lift it to
$t\in\mathbb Z[A(q,r)]$ with $\operatorname{aug}(t)\equiv1\pmod p$.
Then $D(t)=p y$ for some integral vector $y$. The class of $y$ in
$S(r)/\operatorname{im}D(r)$ is killed by both $p$ and $N$, hence is
zero. Write $y=D(z)$ with $z$ integral. Now
$w=t-pz$ lies in $\ker D(r)$ and its augmentation is congruent to
one modulo $p$, so is a nonzero integer. Dividing $w$ by that
augmentation gives the required rational affine equalizer.

Suppose $S$ is generated by $s_j\in S(r_j)$ for $1\leq j\leq m$.
For each $j$ choose $n_j>0$ with $n_js_j\in\operatorname{im}D(r_j)$.
The product $N=\prod_jn_j$ then annihilates the saturation quotient
at every object, by naturality and generation. A subfunctor of a
Noetherian finite sum of representables is finitely generated, which
proves the categorical sufficient condition.

For the group-theoretic sufficient condition, take the filtered union
over finite strong subsets of $M$. The resulting map is
\[
 D_\infty:\mathbb Z[\operatorname{Emb}(q,M)]\longrightarrow
             \mathbb Z[\operatorname{Emb}(p,M)]^k.
\]
Its saturated image is a submodule of the indicated Noetherian
permutation module, hence has finitely many $\mathbb Z[G]$-generators.
An integer multiple of each generator belongs to
$\operatorname{im}D_\infty$, so a common positive integer annihilates
the saturation quotient. The preceding modular lifting argument
works for $D_\infty$. All finitely many embeddings and coefficients
which occur in the resulting rational vector lie in one finite
strong subset of $M$, so it is a rational affine equalizer in $A$.
The all-family affine criterion finishes the proof.
\end{proof}

\paragraph{Remaining total-category gate in C.4.}
The counterexamples above have noncontractible base and hence do not
decide whether all-prime exactness implies rational exactness when
$A$ is additionally assumed total. The arithmetic conclusions of
Proposition~\ref{C:rational-gate} still apply to that restricted
question. Theorem~\ref{C4:contractible-kan-ramsey} excludes a tempting
candidate family: replacing $K(\mathbb Q,r)$ by a contractible Kan
complex in $\Delta_{\mathrm{inj}}/K$ always gives an integrally exact
Ramsey category. The stronger
Theorem~\ref{C:kan-filtered-totalization} excludes every filtered
diagram of connected Kan simplex categories with weakly contractible
colimit base, allowing arbitrary transition functors and arbitrary
coefficient groups. Such a colimit is always integrally exact.
Theorem~\ref{C4:acyclic-cocycle-colimits} also rules out rationally
acyclic filtered colimits of rational-vector coefficient cocycle
categories, without requiring contractibility as a hypothesis.
More generally, Theorem~\ref{C4:exact-base-kan-assembly} excludes
every total Grothendieck assembly of connected Kan simplex
categories over an integrally exact base, even when the base is
not filtered and the transition functors are arbitrary.
The coefficient-pushforward special case in
Theorem~\ref{C4:grothendieck-exclusion} is total precisely when
its coefficient diagram is locally zero.
No total counterexample or proof of the restricted implication is
asserted here.

\begin{corollary}[Proved: a compact torsor witnessing the all-prime failure]
\label{C4:compact-obstruction}
For $A=\mathcal A_2(\mathbb Q)$ there is a diagram $C:A\to\Ab$
that is finitely generated at each object and satisfies
\[
 T_1(C)\cong\mathbb Z,
 \qquad
 \lim_{A^{\mathrm{op}}}^{1} C^\vee\cong\mathbb R/\mathbb Z.
\]
Every profinite abelian diagram on $A^{\mathrm{op}}$ has vanishing
first derived limit, while the displayed compact diagram does not.
A nonzero element of $\mathbb R/\mathbb Z$ gives an objectwise
nonempty compact torsor diagram with empty inverse limit.
\end{corollary}
\begin{proof}
Use the six face embeddings $a_0,\ldots,a_5:P\to Q$ in
Theorem~\ref{C4:counterexample}, and set
\[
 D:\mathbb Z[A(Q,-)]\longrightarrow\mathbb Z[A(P,-)]^5,
 \qquad D(t)=(ta_i-ta_0)_{1\leq i\leq5},
 \qquad C=\operatorname{coker}D.
\]
The rational primitive functional in that proof shows
$\operatorname{aug}\ker D(X)=0$ at every object $X$.
Writing $J=\operatorname{im}D$, right exactness gives
$\operatorname{colim}J\cong\mathbb Z$. Its map to the colimit
$\mathbb Z^5$ of the target representable sum is zero. That sum
is projective, so the long exact sequence of
$0\to J\to\mathbb Z[A(P,-)]^5\to C\to0$ gives
$T_1(C)\cong\mathbb Z$. This argument does not require $D$ itself
to be injective. Finite hom-sets make every $C(X)$ finitely generated.
The duality theorem in Part C gives the asserted first derived limit.

For a direct torsor description, fix $0\ne\theta\in\mathbb R/\mathbb Z$.
Inside
$((\mathbb R/\mathbb Z)^{A(P,X)})^5$ take the affine equations
\[
 \sum_{i=1}^5\bigl(z_i(ga_i)-z_i(ga_0)\bigr)=\theta
       \qquad(g\in A(Q,X)).
\]
Their homogeneous solution group is exactly $C(X)^\vee$.
They are solvable: the homomorphism
$t\mapsto\theta\operatorname{aug}(t)$ on
$\mathbb Z[A(Q,X)]$ vanishes on $\ker D(X)$, hence descends to
$J(X)$ and extends to the ambient free abelian group because
$\mathbb R/\mathbb Z$ is injective. Restrictions along embeddings
preserve the equations. A compatible solution would make each
coordinate function $z_i$ constant on embeddings of $P$, by
restriction along those embeddings from its value at $\operatorname{id}_P$.
The left side of every equation would then be zero, contradicting
$\theta\ne0$. Thus the torsor limit is empty. Finally all-prime
exactness and the profinite test theorem imply vanishing on every
profinite abelian diagram, as claimed.
\end{proof}
